\chapter{Integer Transition Systems}
\chapterauthor{Florian Frohn}
\chapterauthor{Nils Lommen}

\section{Introduction}
This chapter describes the syntax and semantics for \emph{Integer Transition Systems (ITSs)}.
% 
Such systems are \emph{Logically Constrained Term Rewrite Systems} (LCTRSs) where the underlying SMT theory is fixed to the theory of integers, and where function symbols must not occur anywhere but at the root position.
% 
Due to the latter restriction, we only consider top-level rewriting.
% 
Intuitively, such programs correspond to \texttt{while}-programs (or, equivalently, tail-recursive programs) operating on integers.

\section{Syntax}

\subsection{Grammar}
\

\begin{tabular}{lll}
its & ::= & \texttt{(format LCTRS) (theory Ints)} fun+ entrypoint rule+ \\
fun & ::= & \texttt{(fun} IDENTIFIER \texttt{(->} int+\texttt{))} \\
int & ::= & \texttt{Int} \\
entrypoint & ::= & \texttt{(entrypoint} IDENTIFIER\texttt{)} \\
lhs & ::= & IDENTIFIER $|$ \texttt{(}IDENTIFIER IDENTIFIER+\texttt{)} \\
rhs & ::= & IDENTIFIER $|$ \texttt{(}IDENTIFIER int-term+\texttt{)} \\
guard & ::= & \texttt{:guard} formula \\
var & ::= & \texttt{:var (}IDENTIFIER type\texttt{)} \\
rule & ::= & \texttt{(rule} lhs rhs guard? var*\texttt{)} \\
\end{tabular}

\bigskip
Here, ``int-term'' refers to a term of type \texttt{Int} from the underlying SMT-theory \texttt{Ints}, and ``formula'' refers to a term of type \texttt{Bool}.
%
The \texttt{:var} statements are required if the types of variables cannot be inferred (e.g., if the only occurrences of \texttt{x,y} are in the literal \texttt{(= x y)}).

\subsection{Restrictions}

The IDENTIFIER in ``entrypoint'' must be one of the previously declared functions.

\subsection{Remarks}

We allow arbitrary terms and formulas from the \texttt{Ints} theory.
%
In particular, this means that we allow:
\begin{itemize}
\item non-linear arithmetic
\item exponentiation
\item quantifiers
\item \texttt{ite} (if-then-else) expressions
\item \texttt{let} expressions
\item variables of type \texttt{Bool}
\item all usual Boolean connectives (conjunction, disjunction, negation, implication, $\ldots$)
\end{itemize}

We do not require left-linearity, so, e.g.,
\begin{center}
  \texttt{(rule (f x x) (g x))}
\end{center}
is a valid rule.

\section{Semantics}

Let $\F$ be the set of all declared function symbols (i.e., we have $\mathtt{f} \in \F$ if the input file contains a line \texttt{(fun f \ldots)}).
%
In the sequel, we write $f(t_1,\ldots,t_n)$ instead of $\mathtt{(f\ t_1\ \ldots\ t_n)}$ for the sake of readability.
%
Moreover, we always assume that terms respect the arity of the respective function symbols.
%
Thus, if we write $f(t_1,\ldots,t_n)$, then we implicitly assume that $f$ has been declared as $$\mathtt{(fun\ f\ (\texttt{->}\ \underbrace{\mathtt{Int\ \ldots\ Int}}_{n \text{ times}}\ Int))}.$$
%
For a rule with left-hand side $\ell$, right-hand side $r$, and guard $\varphi$, we write $\ell \to r\ \guard{\varphi}$.
%
$\Rules$ denotes the set of all rules that are defined in the input file and $\V$ contains all variables that occur in the input file.
%
The unique entrypoint (as defined by the line \texttt{(entrypoint ...)}) is denoted $\mathit{entry}$.
%
A \emph{configuration} is a term $f(c_1,\ldots,c_n)$ where $f \in \F$ and $c_1,\ldots,c_n \in \Z$.
%
$\CC$ is the set of all configurations.
%
A \emph{valuation} is a mapping $\sigma: \V \to \Z$ and the set of all valuations is denoted by $\Sigma$.
%
$\Rules$ induces a rewrite relation on configurations as follows:
%
We have
\[
  f(c_1,\ldots,c_n) \to g(d_1,\ldots,d_m)
\]
if there is a rule $f(x_1\ldots,x_n) \to g(t_1,\ldots,t_m)\ \guard{\varphi} \in \Rules$ and a valuation $\sigma\in\Sigma$ such that $\sigma(x_i) = c_i$ for all $i \in \{1, \ldots, n\}$, $t_i\sigma \equiv_{\mathtt{Ints}} d_i$ for all $i \in \{1, \ldots, m\}$, and $\sigma \models_{\mathtt{Ints}} \varphi$.
%
Here, $\sigma \models_{\mathtt{Ints}} \varphi$ means that there is an extension $\theta \in \mathtt{Ints}$ of $\sigma\in\Sigma$ such that $\theta \models \varphi$ (note that SMT theories are defined as classes of structures), and $t_i\sigma \equiv_{\mathtt{Ints}} d_i$ means that we have $\theta \models t_i\sigma = d_i$ for all $\theta \in \mathtt{Ints}$.
%
An \emph{evaluation} is a sequence
\[
  \mathit{entry}(c_1,\ldots,c_n) \to s_1 \to \ldots \to s_k.
\]
%
An ITS \emph{terminates} if all evaluations are finite.
%
Let
\[
  ||\sigma||_{\infty} = \max \{ |\sigma(x_j)| \mid 1 \leq i \leq n \}
\]
denote the maximum norm of the state $\sigma\in\Sigma$, i.e., the largest absolute initial value among all program variables $x_1,\dots,x_n$ of $\mathit{entry}(x_1,\ldots,x_n)$.
%
Then the \emph{runtime complexity} is defined as follows.
%
\begin{definition}[Runtime Complexity]
  The \emph{Derivation Height} is the function $\dheight\colon\Sigma\rightarrow\N \cup \{\omega\}$ with
  \[
    \dheight(\sigma) = \sup \{k \in \N \mid \exists s \in \CC.\ \mathit{entry}(x_1, \ldots, x_n)\sigma \rightarrow^k s \}.
  \]
  %
  The \emph{runtime complexity} is the function $\rc\colon\N\rightarrow\N \cup \{\omega\}$ with
  \[
    \rc(k) = \sup \{\dheight(\sigma) \mid ||\sigma||_{\infty} \leq k\}.
  \]
\end{definition}
%
\noindent
%
% A function $rb\colon\Sigma\rightarrow\N \cup \{\omega\}$ is called an \emph{upper runtime bound} if $\rc(\sigma) \leq rb(\sigma)$ for all $\sigma \in \Sigma$.
% %
% Similarly, a function $lb\colon\Sigma\rightarrow\N \cup \{\omega\}$ is called a \emph{lower runtime bound} if $\rc(\sigma) \geq lb(\sigma)$ for all $\sigma \in \Sigma$.
%
\subsection{Termination}
%
Regarding \emph{termination}, the tool provides one of three possible outputs:
\[
  \texttt{YES} \mid \texttt{NO} \mid \texttt{MAYBE}
\]
The tool returns \texttt{YES} if the integer transition system is proven to be terminating and \texttt{NO} if it is non-terminating.
%
Otherwise, \texttt{MAYBE} is returned.
%
\subsection{Complexity}
%
Regarding \emph{complexity}, the tool outputs a combined expression for the worst-case lower and upper bounds on the runtime complexity.
%
The output follows the format:
\[
  \texttt{WORST\_CASE(LB, UB)}
\]
where \texttt{LB} represents the lower bound and \texttt{UB} represents the upper bound (see below for a detailed description).
%
The runtime bounds inferred by the tools are compared asymptotically.
%
For the lower bound \texttt{LB}, the possible outputs and their semantics are:
\begin{itemize}
    \item \texttt{?}: No lower bound could be established.
    \item \text{\texttt{Omega(n\textasciicircum}$d$\texttt{)}}: A polynomial lower bound of degree $d \in \N$, i.e., $rc(n) \in \Omega(n^d)$.
      %
      \begin{itemize}
      \item The output \texttt{Omega(n)} is equivalent to \texttt{Omega(n\textasciicircum 1)}.
      \item The output \texttt{Omega(1)} is equivalent to \texttt{?}.
      \end{itemize}
    \item \texttt{NON\_POLY}: A super-polynomial (e.g., exponential) lower bound.
    \item \texttt{INF}: An infinite lower bound (for non-terminating ITSs, or terminating ITSs with unbounded runtime).
\end{itemize}
%
Similarly, for the upper bound \texttt{UB}, the possible outputs and their semantics are:
\begin{itemize}
    \item \texttt{?}: No upper bound could be established.
    \item \text{\texttt{O(n\textasciicircum}$d$\texttt{)}}: A polynomial upper bound of degree $d \in \N$, i.e., $rc(n) \in \mathcal{O}(n^d)$.
      \begin{itemize}
      \item The output \texttt{O(n)} is equivalent to \texttt{O(n\textasciicircum 1)}.
      \item The output \texttt{O(1)} is equivalent to \texttt{O(n\textasciicircum 0)}.
      \end{itemize}
    \item \texttt{FINITE}: A finite upper bound, e.g., exponential or super-exponential.
\end{itemize}
